\section{Controlled Fourier seeds and the depth needed by parity}
\label{sec:controlled-fourier-seeds}

A weak component already present in the unmodified first-layer kernel
can be amplified at small entropy cost. There is another option: change
the first feature law slightly to insert a stronger component. For parity,
this changes the sufficient depth from a degree-dependent bound to the
same order as for constants and dictators. More generally, the Fourier
$\ell^1$ norm controls a useful family of initial components.

Use the normalized Fourier expansion on the Boolean cube,
\[
 f=\sum_{S\subseteq[d]}\widehat f(S)\chi_S,\qquad
 \chi_S(x)=\prod_{i\in S}x_i,\qquad
 A(f)=\sum_S|\widehat f(S)|.
\]
For Boolean $f$, Parseval implies $A(f)\ge1$.

\begin{theorem}[A controlled Fourier-seed bound]
\label{thm:controlled-fourier-seed}
There are universal constants $C_0,C_1$ such that, if
\[
 T_f=\log(d+L+2)+\log A(f),\qquad L\ge C_0T_f,
\]
then
\[
 J_L(f)\le V_*+C_1\frac{T_f^2}{L}.
\]
In particular, a depth at least
$C[\log(d+2)+\log A(f)]^2$ puts $f$ below a fixed universal absolute
critical-rate budget.
\end{theorem}

\begin{corollary}[Parity and juntas]
\label{cor:controlled-parity-junta-depth}
Every parity, irrespective of its degree, attains a fixed universal
critical-rate budget at depth $O(\log^2(d+2))$.
The universal lower bound in Theorem~\ref{thm:critical-entropy-transport}
makes $\Theta(\log^2(d+2))$ the necessary and sufficient depth scale
for attaining some fixed universal budget, as $d$ grows.
Every function depending on at most $k$ coordinates attains such a
budget at depth $O((k+\log(d+2))^2)$.
These are rate statements. Corollary~\ref{cor:critical-fourier-learning}
uses them to obtain polynomial sample bounds, already at the stated
logarithmic minimum learner depth.
\end{corollary}

\subsection{A common polynomial-cost family of parity laws}

\begin{lemma}[Parity mean and constant diagonal]
\label{lem:controlled-parity-law}
For every $S\subseteq[d]$, there is a first-feature law $\nu_S$ with
\[
 D(\nu_S\Vert\mathsf M_{Q_0})\le D_*:=C(d+1)^3\log^2(d+2),
 \qquad \mathbb E_{\nu_S}H=c_S\chi_S,
 \qquad c_S\ge\frac1{8(d+1)}.
\]
Its second-moment diagonal is constant and belongs to $[F(1),1]$, where
$F(q)=\mathbb E\tanh^2(\sqrt qZ)$ as above. Increase the universal $C$
if necessary so that $D_*\ge1$.
\end{lemma}

\begin{proof}
First construct a law with a signed mean margin, then symmetrize it.
Let $k=|S|\ge1$ and $u(x)=\sum_{i\in S}x_i$.
At its possible values $u_j=-k+2j$, the parity equals
$f_j=(-1)^{k-j}$. Thus
\[
 \chi_S(x)=c_0+\sum_{j=1}^k c_j
       \operatorname{sign}(u(x)-(-k+2j-1)),
 \qquad |c_0|\le1,\quad c_j=(f_j-f_{j-1})/2\in\{-1,1\}.
\]
Put $M=4+\log(64(k+1))$.
Shift a first-row Gaussian weight vector to have mean coefficient $M$
on each coordinate in $S$, zero on the other inputs, and bias
$-M(-k+2j-1)$. Its preactivation mean has the required sign and
absolute value at least $M$ at every cube input. Its variance remains
$q_0=(d+1)/d\le2$. For a centered Gaussian $G$ of variance at most two,
\[
 1-\mathbb E\tanh(t+G)\le2e^{-2t+4}\qquad(t\ge0).
\]
The expected feature $v_j$ therefore approximates its ideal step to
maximum error $1/[16(k+1)]$.
A constant bias shift supplies a constant mean $c_{\rm const}\ge1/2$.

Take $\lambda=1/[4(k+1)]$. Mix the appropriately signed detector
laws with probabilities $\lambda|c_j|$, and the appropriately signed
constant law with probability $\lambda|c_0|/c_{\rm const}$.
Their sum is at most $\lambda(k+2)\le1$; use the unmodified law for
the remaining probability. Its feature mean is zero. The resulting
mean $m_S$ obeys
\[
 \chi_S(x)m_S(x)\ge\lambda\left(1-\frac{k}{16(k+1)}\right)
                  \ge\frac1{8(k+1)}
 \qquad\text{for every }x.
\]
Every detector shift has Gaussian relative entropy at most
\[
 \frac d2 M^2\{k+(-k+2j-1)^2\}
 \le C(d+1)(k+1)^2\log^2(k+2).
\]
Convexity and data processing give the common bound $D_*$ for this
feature law. The case $k=0$ uses only a positive constant bias shift,
with constant positive mean and entropy $O(d)$.

For $b\in\{-1,1\}^d$, transform features by
\[
 (T_bH)(x)=\chi_S(b)H(b_1x_1,\ldots,b_dx_d).
\]
Input weight sign changes and a simultaneous sign change of all weights
implement this transformation. It preserves the raw first Gaussian
law. Average the transformed controlled laws over $b$.
Their entropies are unchanged individually, so convexity bounds the
mixture's entropy by $D_*$. Its mean is exactly
\[
 \mathbb E_b\chi_S(b)m_S(bx)
   =c_S\chi_S(x),\qquad
 c_S=\mathbb E_x\chi_S(x)m_S(x)\ge\frac1{8(k+1)}.
\]
Transitivity of cube translations makes its diagonal second moment
constant. Each original component was a Gaussian mean shift with
unchanged covariance. For centered Gaussian $G$, $|G+t|$ stochastically
dominates $|G|$, so
$\mathbb E\tanh^2(G+t)\ge\mathbb E\tanh^2(G)$.
Every diagonal before and after symmetrization is consequently at least
$F(q_0)\ge F(1)$ and at most one.
Only the first feature law has been symmetrized; no convexity of entire
covariance trajectories is asserted.
\end{proof}

\subsection{From a Fourier mixture to a small initial entropy cost}

\begin{proof}[Proof of Theorem~\ref{thm:controlled-fourier-seed}]
Negating a whole feature preserves the raw first law and negates its
mean. Choose the corresponding sign of each law in
Lemma~\ref{lem:controlled-parity-law}, and set
\[
 Z_f=\sum_S\frac{|\widehat f(S)|}{c_S}
       \le8(d+1)A(f).
\]
Mix these laws with probabilities
$|\widehat f(S)|/(c_SZ_f)$. The resulting law $\nu_f$ has entropy
at most $D_*$, mean exactly $f/Z_f$, and second moment
\[
 Q(\nu_f)\succeq Z_f^{-2}ff^T.
\]
Its diagonal is constant in $[F(1),1]$. The component diagonals may
have different values, but each is constant over inputs. Convexity
introduces no factor for the number of Fourier components.

Now replace the original first law by
\[
 \widetilde\nu_f=(1-\theta)\mathsf M_{Q_0}+\theta\nu_f,
 \qquad \theta=(D_*L)^{-1}.
\]
Its entropy is at most $1/L$, its diagonal remains constant in
$[F(1),1]$, and
\[
 Q(\widetilde\nu_f)\succeq vff^T,
 \qquad v=\frac1{D_*LZ_f^2}.
\]
This uses the second moment of the mixture; it does not square the
mixing probability. Hence
\[
 S_v:=1+\log(1/v)\le C[\log(d+L+2)+\log A(f)]=CT_f.
\]

We spell out the initial-condition adjustment in the existing steering
argument. Lemma~\ref{lem:critical-free-cooling} is stated for an
initial diagonal in $[F(1),F(2)]$, but its proof works on $[F(1),1]$.
Indeed $F(q)<q$, $(q-F(q))/q^2$ has a positive lower bound on $(0,1]$,
and
\[
 R(q)=\frac{q[1-F(q)]^2}{F(q)}>0,
 \qquad \log R(q)=O(q^2)\quad(q\downarrow0).
\]
Compactness extends the required estimates to the enlarged interval.
The same telescoping proof reaches diagonal $\delta$ in $O(1/\delta)$
free layers while retaining a component at least $c\delta v$.

The proof of Theorem~\ref{thm:critical-target-seed-upper} then applies
with $S_v$ in place of $S_f$ and with the first-law entropy $1/L$
added. For clarity, set $\delta=B_1S_v/L$ and
$p=B_2(S_v+\log(2L))/L$ with the sufficiently large universal constants
used there. Prescribed rank-one growth costs $O(p+S_v\delta)$ and
uses $O(S_v/\delta)$ layers. Exact holding in the unused layers
costs $O(L\delta^2)$. Purification costs $O(p\log(2L))$, leaves
residual $O(\delta L^{-4})$, and the scalar path costs at most
$V(p)+p\le V_*+p$, with $O(L^{-4})$ additional readout cost.
All stage lengths fit when $L\ge CS_v$, and holding gives the exact
specified depth. Adding the initial entropy yields
\[
 J_L(f)\le V_*+\frac1L+
 C\frac{S_v^2+\log^2(2L)}L
 \le V_*+C_1\frac{T_f^2}L.
\]
The strict-margin rate equals its closed-margin variational infimum
by Corollary~\ref{cor:critical-margin-rates}, so this construction
has precisely the prior-rate interpretation in the theorem.

Finally let $B=\log(d+2)+\log A(f)\ge1$.
One has $T_f\le B+\log(L+1)+C$. For every $L\ge CB^2$, increasing
$C$ if necessary makes $T_f/L$ satisfy the depth requirement and
bounds $T_f^2/L$ by a universal constant. This proves the fixed-budget
assertion.
\end{proof}

\begin{proof}[Proof of Corollary~\ref{cor:controlled-parity-junta-depth}]
A parity has $A(f)=1$. For a $k$-junta the Fourier support has at most
$2^k$ elements, so Parseval and Cauchy--Schwarz give
$A(f)\le2^{k/2}$. Substitute these bounds into the theorem.
The matching necessary depth for any fixed rate budget follows from
Theorem~\ref{thm:critical-entropy-transport}.
\end{proof}

The controlled mixture explains why a small eigenvalue in the original
first-layer kernel cannot by itself give a large critical-rate lower
bound. The mixture inserts a stronger component at entropy $1/L$.
The theorem gives no lower bound ordering the competing extensions of
a training sample; the next section isolates that remaining question.
